\documentclass[aps,prl,reprint,10pt,amsmath,amssymb,superscriptaddress,a4paper]{revtex4-2}
\usepackage{graphicx} % Graphics package
\usepackage{dcolumn} % Double column package
\usepackage{amsmath,amsfonts,amsthm} % Math packages
\usepackage[margin=2cm]{geometry} % Sets 2cm margins
\usepackage{datetime} % Package for automatic date & time
\usepackage{subfig}
\usepackage[hidelinks]{hyperref}

\begin{document}

\title{Chaotic Circuit}

\author{T. Nguyen (z5416116)}
\affiliation{UNSW PHYS3117 - Physics Laboratory}
\date{\currenttime~\today}

\begin{abstract}
l
\end{abstract}

\maketitle

\section{Introduction}

\begin{quote}
\small
"Given for one instant an intelligence which could comprehend all the
forces by which nature is animated and the respective
situation of the beings who compose it-an intelligence
sufficiently vast to submit these data to analysis-it
would embrace in the same formula the movements of
the greatest bodies of the universe and those of the
lightest atom; for it, nothing would be uncertain and
the future, as the past, would be present to its eyes."
\begin{flushright}
- Pierre Simon Laplace, \textit{A philosophical essay on probabilities, pg. 4} \citep{laplacedemon}
\end{flushright}
\end{quote}

\vspace{1em}

The topic of determinism has been discussed by philiosophers and physicists alike for millenia and will continue 
for the next thousands of years \citep{sep-freedom-ancient}. In a philosophical sense, determinism is more about 
free will and moral responsibility of human beings for their actions, whereas the interpretation in physical 
sciences, relates to predictions given the present state of systems. It is outside the scope of this report to, 
firstly, argue the necessity for determinism in governing this universe and, secondly, discuss the interplay of free will 
and determinism, we will leave this to contemporary philiosophers and budding online intellectuals. 

In the realm of chaos theory, I will assume that determinism is the correct interpretation of our universe, that is, 
to assume that the demon is real and computationally possible. The definition of determinism I will 
be abiding by is one given by Hoefer's paper on "Causal Determinism" (2003), \citep{sep-determinism-causal}, where 
determinism is loosely encompassed by the definition: 

\begin{quote}
    Determinism is true of the world if and only if, given a specified way things are at a time t, the way things go 
    thereafter is fixed as a matter of natural law.
\end{quote}

In other words, the world this report resides in, is one governed solely by Newton's classical laws of physics. I will 
ignore the challenges to determinism that quantum mechanics brings \citep{earman} and instead focus on the certainty of states given 
their initial conditions. The study of chaos is, counter-intuitively, not probabilistic. Chaotic systems do not have 
different outcomes if a perfect process is repeated. Instead, the different outcomes only appear to us humans because 
of our imperfect knowledge of the parameters of the initial state of the system. Therefore, to an all knowing demon 
who is able to perceive the parameters of a system to infinite figures, the future and the past, are one and the same.

The language of chaos is often, understandably so, misused as a synonym for randomness. This comes from their shared 
characteristic of unpredictability however the root cause of this unpredictability is not the same. Chaotic systems are 
inherently unpredictable due to our lack of mathematical precision and the system's incredible sensitivity to inital 
conditions. The time evolution of these systems may or may not exhibit patterns to us which also propagates the notion of 
randomness however they are not equivalent. It is important to clear this language as throughout this report, I may describe 
trajectories or shapes as disordered and lacking form but this does not mean it is not governed by deterministic physical law.

The conditions for chaos are well understood, being formulated at the time of the discovery of chaotic behaviour by Edward Lorenz 
in 1963 \citep{lorenz}. The discovery was serendpitious and as the story goes, Lorenz was trying to model weather patterns using 
the state space of a 12 component model. Each component possessed a differential equation relating to a real world atmospheric 
parameter such as temperature, pressure etc. He tried to solve this model numerically however found that his model was nonperiodic 
whereas most studied solutions were periodic.

The famous story in his career is that one day, he wanted to examine one of the solutions in greater detail. He had inputted in 
12 numbers that the simulation had printed out, effectively restarting the process from about half way and found that his two runs
had completely different outputs despite their shared initial conditions. Later, he would attribute the truncated output to be the 
source of this divergence \citep{lorenzbio}.

The computer he used for the simulation stored the output numbers to six decimal places however when he printed them out to restart 
the process, he only printed up to three decimal places to save paper. This one part in ten thousand difference had caused 
the simulation to wildly alter its course from its original trajectory \citep{lorenzstory}. The first condition of chaotic motion is 
then the system must exhibit sensitive dependence on its inital conditions. 

The story then goes that while visiting his colleague, Barry Saltzman, he found that in Barry's work, there was a three component subset
that displayed aperiodicity, or chaotic motion. This three component model became the famous Lorenz model for 
which chaotic motion can arise \citep{sparrow}. The beauty of the Lorenz equations came from their simplicity, only requiring three variables
for the properties of chaos to appear. Through his work and shown in his paper in 1963, another condition for the aperiodicity is the 
equations must be non-linear.

Remarkably, the development of chaos theory since its discovery in 1963, continues on the work of Newton in his final pursuit 
for a solution in the three body problem. Coincidentally, the three body problem has conditions 
that align with the ones set out by chaos theory. Newton successfully solved the two body problem when he introduced to the world his idea 
of gravitation. In that scenario, he only needed to worry about the motion of two celestial objects and therefore, the problem could be 
solved analytically and the solution exhibited periodicity and stability always. However, the inclusion of a third gravitational body would 
introduce a non-linear term into the equations as well as increasing the variable count to three, satisfying all conditions for aperiodicity 
and therefore the solution to Newton's three body problem is chaotic \citep{galaviz}.

The word 'chaos' was first coined by Jim Yorke to describe types of non-linear dynamical systems however the term has transcended the technical 
community. The idea of chaos and its child, 'The Butterfly Effect', has become apart of everyday language and integrated itself seamlessly in pop 
culture. It is common to hear in economics and political discussions, the term 'bifurcation' in relation to causal relationships, especially one 
where predictability of outcomes is at its lowest. There is a Wikipedia article dedicated to the Butterfly Effect in popular culture as the metaphor 
of a butterfly flapping its wings to cause a tornado in Texas (presented by Lorenz in a 1972 talk) resonates with the general public and their 
perception of the world around them. The imagery of the smallest detail influencing the world's most significant events is one of the most powerful 
narratives to ever been told. But the best kind of stories are the ones that did actually happen and its our marvel of the possibility turning 
into reality is the reason why the butterfly effect and chaos theory really has a special place in our society today.

\section{Theory}

The central characteristic of chaos is that the system does not repeat its past behaviour, not even approximately. Despite their lack of regularity, 
the aperiodicity still follow the deterministic equations derived from Newton's laws.

\subsection{Conditions for chaos}

As mentioned before, Edward Lorenz had come up with the two conditions for chaotic motion. For a system to be chaotic, it must:

\begin{enumerate}
    \item be described by at least three components that are independent dynamical variables
    \item contain a non-linear term in its equations of motions that couples several of the variables
\end{enumerate}

The sensitivity to initial conditions is a consequence of the second condition. THe non-linear term grows errors in predictions exponentially in time so 
any small shifts from its original state will cause that state to become essentially unrecognisable in evolution. It is this consequence that will drive 
the motivation for investigation in chaos theory. Another non-trivial consequence of the second condition is that differential equations containing non-linear 
terms are impossible to solve analytically and so solutions for equations of motion will have to be done numerically.

\subsection{Phase space}

To study dynamical systems, often phase spaces are used to describe states and their evolution. The phase space is a mathematical space with orthogonal coordinate 
directions representing each of the variables required to define the state of the system. Commonly, the axis of the phase space will come in position, velocity pairs 
or analogues. For each solution on a phase space diagram, there will be a corresponding phase trajectory. 

There is an extremely non-trivial and significant property 
of these trajectories, being they are non-crossing. Two trajectories, no matter how close, will never cross each other. This is a consequence of determinism, once a 
state of the system is known, its time evolution into the future and in the past are known as well. A state in two trajectories would render the system indeterminate 
and therefore contradict the uniqueness of the determinism on which the motion is formulated upon. 

Phase spaces of conservative (constant energy) systems preserve the areas in the space. All the points in a given area of phase space will move in such a way to preserve 
the area occupied by these same points at some later time. If the system does not conserve energy, the phase area or volume will contract. This is known as a dissipative 
dynamical system and they introduce a mathematical object known as an attractor. The point is defined as the end state, or fully time evolved state, for all finite states 
that describes the system.

In order to simplify complicated phase space diagrams, Poincare sections are used. They are constructed by viewing the phase space diagram stroboscopically in such a way that 
the motion is observed periodically. Therefore, the Poincare section will provide information about the ratio of the strobing and natural frequency of the dynamical motion. If 
the ratio is rational such that the natural frequency $\omega_0$ is equal to $\frac{p}{q}\omega_s$ where $\omega_s$ is the strobing frequency and $p$ and $q$ are integers, then 
there are $q$ points on the diagram. If the ratio is irrational, then no points are repeated on the plot and so a filled in line of the trajectory will appear.

\section{Experimental Setup}

\section{Results and Discussion}

\section{Conclusion}

\section{Acknowledgements}

\bibliographystyle{apsrev4-2}
\bibliography{references}

\end{document}